
\reportelement{Введение}

Цель работы — развернуть приложение FastAPI из ЛР №4 в локальном кластере Kubernetes. Работа выполнена в Windows PowerShell с Docker Desktop и Minikube 1.39.0. Версия Kubernetes — 1.37.0. Необходимо проверить доступность приложения, запись заметок в PostgreSQL и кэширование погоды в Redis.

\begin{table}[H]
  \caption{Персональные параметры}\label{tab:personal}
  \centering
  \begin{tabularx}{\textwidth}{|p{0.42\textwidth}|X|}
    \hline
    Параметр & Значение\\ \hline
    Номер ИСУ & \StudentISU\\ \hline
    Префикс ресурсов & \texttt{yamaleeva}\\ \hline
    Имя базы данных & \texttt{yamaleeva83}\\ \hline
    Секретный ключ & \texttt{secret83yamaleeva}\\ \hline
    Порт NodePort & \texttt{30083}\\ \hline
    Значение \texttt{APP\_OWNER} & \StudentName\\ \hline
  \end{tabularx}
\end{table}

Ветка работы: \url{\RepositoryURL}.

\chapter{Ход выполнения}

\section{Подготовка образа}

Приложение собирается из \texttt{Dockerfile} в каталоге проекта.
Образ загружается в профиль \texttt{lab5}. Политика \texttt{Never}
указывает использовать этот локальный образ.

\reportscreenshot{01-image.png}{Узел в состоянии Ready и образ приложения в Minikube}{fig:image}

\section{Развёртывание PostgreSQL и Redis}

Для постоянного хранения данных PostgreSQL создан PVC на 1 ГиБ
с режимом доступа \texttt{ReadWriteOnce}.

\manifestshot[0.5]{postgres-pvc}{Манифест PVC для PostgreSQL}

\clearpage
PostgreSQL запускается через Deployment и доступен внутри кластера
через сервис \texttt{ClusterIP} на порту 5432.

\manifestshot[0.79]{postgres-deployment}{Манифест Deployment PostgreSQL}

\manifestshot[0.33]{postgres-service}{Манифест сервиса ClusterIP для PostgreSQL}

\clearpage
Redis развёрнут через Deployment. Сервис \texttt{ClusterIP}
принимает внутренние запросы на порту 6379.

\manifestshot[0.48]{redis-deployment}{Манифест Deployment Redis}

\manifestshot[0.38]{redis-service}{Манифест сервиса ClusterIP для Redis}

\clearpage
Поды PostgreSQL и Redis перешли в состояние \texttt{1/1 Running}.
Том PostgreSQL получил статус \texttt{Bound}.

\reportscreenshot{02-resources.png}{Поды PostgreSQL и Redis, сервисы и привязанный PVC}{fig:resources}

\section{Развёртывание приложения FastAPI}

Для ключа приложения и пароля PostgreSQL используется Secret.
В репозиторий добавлен файл \texttt{app-secret.example.yaml}
со значениями-заглушками.

\manifestshot[0.75]{app-secret.example}{Файл шаблона Secret}

\clearpage
Из ConfigMap приложение получает ФИО, параметры PostgreSQL,
адрес Redis и время хранения записи в кэше.

\manifestshot[0.59]{app-configmap}{Манифест ConfigMap приложения}

Значение \texttt{APP\_OWNER} проверено непосредственно в ConfigMap.
Оно совпало с ФИО в ответе API.

\reportscreenshot{06-owner.png}{Значение APP\_OWNER в ConfigMap}{fig:owner}

\clearpage
Deployment FastAPI запускает одну реплику. Init-контейнер ждёт
готовности PostgreSQL. Настройки берутся из ConfigMap и Secret.

\manifestshot{app-deployment}{Манифест Deployment приложения FastAPI}

\clearpage
Приложение слушает порт 8083. Сервис типа NodePort публикует порт 30083.

\manifestshot[0.42]{app-service}{Манифест сервиса NodePort}

После загрузки обновлённого образа FastAPI запустился.
В выводе сервиса видны порты \texttt{8083:30083/TCP}.

\reportscreenshot{02-app.png}{Запуск FastAPI и сервис NodePort}{fig:app}

\clearpage
\section{Проверка ресурсов и API}

Команда \texttt{kubectl get all} показала поды, сервисы, Deployment
и ReplicaSet в пространстве имён \texttt{default}.
Все три пода находились в состоянии \texttt{1/1 Running}.
Журнал пода FastAPI показывает ответы HTTP 200 до обновления проб.

\reportscreenshot{03-resources.png}{Ресурсы пространства имён default}{fig:all}
\reportscreenshot{03-logs.png}{Журнал контейнера FastAPI}{fig:logs}

Команда \texttt{minikube service} открыла доступ к сервису
на локальном адресе \url{http://127.0.0.1:7127}.
В другом окне PowerShell через этот адрес выполнены запросы
\texttt{GET /} и \texttt{GET /health}.

\reportscreenshot{03-service.png}{Локальный адрес сервиса Minikube}{fig:service}
\reportscreenshot{03-http.png}{Ответы GET / и GET /health через сервис Minikube}{fig:http}

В ответе \texttt{GET /} отображаются приветствие, имя пода
(\texttt{hostname}) и время сервера (\texttt{server\_time}).
В \texttt{GET /health} получены \texttt{status: healthy},
\texttt{postgres: true} и \texttt{redis: true}. Поле \texttt{owner}
содержит ФИО, заданное переменной \texttt{APP\_OWNER}.

Для проверки кэша дважды запрошена погода для Казани.
Первый ответ содержал \texttt{source: api}, второй через несколько
секунд содержал \texttt{source: cache}. TTL кэша составляет 120 секунд.

\reportscreenshot{04-weather.png}{Первый и повторный запросы погоды}{fig:weather}

Запрос \texttt{POST /notes} сохранил заметку
\texttt{LR5 PostgreSQL test} с идентификатором 1.
Запрос \texttt{GET /notes} вернул ту же заметку с \texttt{id: 1}.

\reportscreenshot{05-notes.png}{Добавление и чтение заметки}{fig:notes}

\clearpage
\section{Пробы работоспособности}

Обе HTTP-пробы из манифеста на рисунке~\ref{fig:notepad-app-deployment}
обращаются к \texttt{/health} на порту 8083.
Для readiness заданы задержка 5 секунд и период 10 секунд.
Для liveness заданы задержка 15 секунд и период 20 секунд.

Readiness определяет, может ли под принимать запросы через Service.
При неудачной проверке Kubernetes временно исключает под из получателей
трафика. Liveness определяет, нужно ли перезапустить контейнер.
При трёх последовательных сбоях liveness kubelet перезапускает
контейнер. В этой конфигурации \texttt{/health}
проверяет PostgreSQL и Redis, поэтому продолжительная недоступность
одной из этих служб может вызвать перезапуск приложения.

После применения обновлённого Deployment новый под получил
\texttt{Ready=True}. В выводе \texttt{kubectl describe pod}
обе пробы показаны с путём \texttt{/health}.

\reportscreenshot{06-probes.png}{Обновление Deployment и проверка проб пода}{fig:probes}

\clearpage
\section{Дополнительное задание. Ingress}

В Minikube используется аддон Ingress с контроллером NGINX.
Правило для хоста \texttt{yamaleeva.local} направляет запросы на
сервис \texttt{yamaleeva-app} через порт 8083.

\manifestshot[0.58]{app-ingress}{Манифест Ingress для yamaleeva.local}

\reportscreenshot{07-ingress.png}{Правило Ingress и проброс порта контроллера}{fig:ingress}

Из Windows доступ к контроллеру проверен через локальный порт 18084.
Параметр \texttt{--resolve} указал curl использовать адрес 127.0.0.1
для хоста \texttt{yamaleeva.local}. Запрос \texttt{/health} через
Ingress вернул \texttt{healthy}. Имя в поле \texttt{hostname}
совпало с именем пода после обновления Deployment.

\reportscreenshot{07-health.png}{Ответ FastAPI на запрос через Ingress}{fig:ingress-health}

\reportelement{Заключение}

В Minikube развёрнуты FastAPI, PostgreSQL и Redis. Все три пода
перешли в состояние \texttt{Running}, а PVC PostgreSQL получил статус
\texttt{Bound}. Через API проверены сохранение заметки и чтение погоды
из кэша. Для пода FastAPI подтверждены состояние \texttt{Ready}
и наличие обеих проб на \texttt{/health}. Дополнительно проверен доступ к приложению
через Ingress для хоста \texttt{yamaleeva.local}.

